\section{Conclusion and Future Outlook}
\label{sec:conclusion}

In this work, we have challenged the traditional compromise-driven assumptions of static model merging and the batch-dependency constraints of existing dynamic ensembling methods.
We proposed \textbf{EpiMerge} (\textbf{Epigenetic Weight Masking}), a biologically-inspired model-merging framework that enables true sample-specific dynamic ensembling in parallel.
By mimicking cellular molecular regulatory layers, EpiMerge translates the input sample itself into an environmental stimulus, generating low-rank row-wise and column-wise gating masks that dynamically and reversibly scale pre-trained expert parameters.
Crucially, using vectorized tensor contractions (\texttt{torch.einsum}), EpiMerge executes sample-wise weight reconstruction and parallel inference in a single GPU pass, bypassing batch-averaging ensembling shortcuts.

Our empirical evaluation across MNIST, FashionMNIST, CIFAR-10, and SVHN on a Vision Transformer backbone demonstrated that EpiMerge is completely immune to temporal non-I.I.D. task shifts (Bursty) and extreme small-batch noise ($B=2$), outperforming standard static and online test-time adaptive baselines.
By maintaining perfect sample-wise specificity and independence, EpiMerge delivers highly predictable and stable multi-task inference.

\subsection{Future Outlook}
From the perspective of a visionary future, we believe that epigenetic weight masking represents a foundational paradigm shift for highly adaptive neural systems. We identify several exciting and open avenues of future research:

\begin{enumerate}
    \item \textbf{Scaling to Hundreds of Tasks:} While this work evaluated four expert models, the parameter-efficient, low-rank design of Epigenetic Reader Heads makes them ideally suited for large-scale multi-task ensembling. Future work will explore merging dozens or hundreds of specialized experts, where the low-rank latent space can act as a continuous coordinate map of specialized machine learning domains.
    \item \textbf{Epigenetic Large Language Models:} The mathematical formulation of EpiMerge can be seamlessly integrated into transformer layers of Large Language Models (LLMs) and multi-modal foundational architectures. By dynamically scaling the query, key, value, and feed-forward projections at the token level, an LLM could dynamically adapt its underlying weights to express different functional personas, styles, or domain expertises per token or sentence.
    \item \textbf{Evolutionary and Self-Supervised Adaptation:} Instead of relying on a tiny supervised offline calibration dataset, future iterations can explore using reinforcement learning or open-ended evolutionary algorithms to adapt the Epigenetic Reader Heads on-the-fly. This could enable models to dynamically discover optimal parameter scaling paths in completely novel test environments.
    \item \textbf{Synaptic Resiliency and Lifelong Learning:} By treating the pre-trained weights as a permanent genetic memory and the epigenetic masks as a plastic, fluid adaptation layer, neural networks can achieve robust lifelong learning. New skills can be acquired by training only the tiny epigenetic reader heads, while the static base weights remain completely untouched, providing a mathematically guaranteed defense against catastrophic forgetting.
    \item \textbf{Lightweight Sensory Extractors and Edge Deployment:} Although the proposed dual-stream sensory extractor copy provides highly robust semantic guiding signals, it doubles the static parameter footprint and triples wall-clock latency (from 9.12ms to 27.34ms at $B=64$). While our proposed \textbf{EpiMerge-Active} variant elegantly slashes this overhead to exactly $1.0\times$ parameters and $1.0\times$ latency by extracting representations from early active layers, it incurs a substantial performance penalty, dropping from $39.22\%$ to $36.70\%$. This significant system-accuracy trade-off highlights a major deployment challenge. A highly promising and critical next step for practical edge deployment is the development of ultra-lightweight sensory extractors. Future work will investigate utilizing extremely compact distilled student models (e.g., a tiny MobileNet or a 1-layer MLP) to generate guiding signals, or implementing hardware-efficient pointer weight-sharing between active and sensory backbones to minimize GPU memory bandwidth bottlenecks.
\end{enumerate}

By merging the insights of molecular biology with the mathematical beauty of deep parameter manifolds, EpiMerge provides a stepping stone toward a new class of self-organizing, adaptive, and resilient artificial intelligence systems.
